% Shared exercise chunk: l4-persistent-memory
\begin{exercisechunk}
\item \textbf{Forgetting errors and retaining information.}
\label[exercise]{ex:persistent-memory}
Use the reset mixture and persistent-state model from
\cref{sec:selective-memory}.
\begin{enumerate}[(a),beginpenalty=10000]
\item Prove the contraction bound for the reset mixture \(K_r\).
Construct a process whose next-token law uses only the last token
and whose initial bit can be recovered from every later token.
\item In \cref{eq:persistent-working-memory}, suppose the conditional
working-state kernels contract uniformly with factor \(\rho<1\).
Explain why contraction applies when comparing equal \(M\), while
different persistent \(M\)'s can remain distinguishable indefinitely.
\end{enumerate}

\end{exercisechunk}
